% Generated draft: requires mathematical and source review.
\documentclass[11pt]{article}
\usepackage[margin=25mm]{geometry}
\usepackage{fontspec}
\setmainfont{Latin Modern Roman}
\usepackage{amsmath,amssymb,mathtools,unicode-math}
\setmathfont{Latin Modern Math}
\usepackage{enumitem,needspace,xurl,hyperref}
\setlength{\parindent}{0pt}
\setlength{\parskip}{5pt}
\setlength{\emergencystretch}{3em}
\newcommand{\N}{\mathbb{N}}
\newcommand{\Z}{\mathbb{Z}}
\newcommand{\Q}{\mathbb{Q}}
\newcommand{\R}{\mathbb{R}}
\newcommand{\C}{\mathbb{C}}
\newcommand{\F}{\mathbb{F}}
\newcommand{\A}{\mathbb{A}}
\newcommand{\PP}{\mathbb{P}}
\newcommand{\E}{\mathbb{E}}
\newcommand{\eps}{\varepsilon}
\newcommand{\dd}{\,\mathrm{d}}
\begin{document}
\section*{UnsolvedMath ID: 700942}
\textbf{Source problem code:} catalogue entry
\subsection{Definitions and mathematical statement}
Let \(G\) be a topological group with identity \(e\), whose topology is not discrete. An ultrafilter on the underlying set of \(G\) means a maximal proper filter of subsets of \(G\). Such an ultrafilter \(\mathcal U\) converges to \(e\) if every neighbourhood of \(e\) belongs to \(\mathcal U\). Put
\[
\mathscr U_e(G)=\{\mathcal U:\mathcal U\text{ is an ultrafilter on }G\text{ and }\mathcal U\to e\}.
\]
Assume that \(\mathscr U_e(G)\) is finite. Must there exist a subgroup \(H\leq G\) such that
\[
H\text{ is open in }G,\qquad |H|\leq\aleph_0,\qquad h^2=e\quad\text{for every }h\in H?
\]
Here countable means finite or countably infinite. The condition on squares expresses exponent \(2\): since \(G\) is non-discrete, an open subgroup cannot be \(\{e\}\). Ultrafilters are understood to include principal ultrafilters; in particular, the principal ultrafilter at \(e\) is included in \(\mathscr U_e(G)\).

The supplied statement does not specify a separation axiom for topological groups. Thus the intended convention, particularly whether Hausdorffness is required, remains unresolved; the formulation above records the stated hypotheses without adding that requirement.

The supplied source labels the general problem open. Its comment dated 2005 reports that the conclusion was proved when \(G\) is countable, citing E. G. Zelenyuk, \emph{Mat. Studii} \textbf{7} (1997), 139--144. This is a supplied historical qualification, not an independently verified result or current literature assessment.
\subsection{Short English statement}
Must every non-discrete topological group with only finitely many ultrafilters converging to its identity contain a countable open subgroup in which every element squares to the identity? The intended separation convention needs clarification. A supplied 2005 comment reports an affirmative result for countable groups.
\subsection{Sources}
\begin{itemize}
\item Local UnsolvedMath registry: permanent ID \(700942\), exact source code \texttt{catalogue entry}, title \emph{Kourovka Notebook Problem 13.24}. The supplied registry status is \texttt{open}; this label has not been independently verified.
\item Attributed to E. G. Zelenyuk, proposed in 1995. Supplied source: \emph{The Kourovka Notebook (Unsolved Problems in Group Theory)}, 13th issue (1995), Problem 13.24; also identified with the 21st issue edition (2026), arXiv:1401.0300v48. The supplied collection description names E. I. Khukhro and V. D. Mazurov as editors.
\item Supplied comment dated 2005: the countable case is reported proved by E. G. Zelenyuk, \emph{Mat. Studii} \textbf{7} (1997), 139--144. No article title or further bibliographic details are supplied. Neither this citation nor the reported result has been independently checked.
\end{itemize}
\end{document}
